% !TeX root = ../../main.tex

PostgreSQL orientiert sich an den SQL Standardkonventionen für \texttt{SQLSTATE} codes.
Daher sind die Rückgabewerte für Transaktionen fünfstellig und bestehen sowohl aus Zahlen, als auch aus Buchstaben.

Für die Arbeit sind die folgenden Fehlercodes relevant.
\begin{itemize}
	\item \texttt{00000}: successful completion
	\item \texttt{40001}: Serialization Failure
	\item \texttt{40P01}: Deadlock detected
	\item \texttt{55P03}: Lock not available
\end{itemize}
Der Fehler \texttt{40001 Serialization Failure} tritt unter anderem bei dem \emph{Serializable Isolation Level} auf.
Dieses Level garantiert, dass gleichzeitig laufende Transaktionen kein Ergebnis erzeugen dürfen, der nicht entstehen
würde, wenn die Transaktionen nacheinander ausgeführt würden.
Falls das passiert, wird für eine Transaktion ein Serialization Failure Fehler zurückgegeben, während eine andere
erfolgreich abschließen kann~\cite{PostgreSQLSerializable}.

Ein Deadlock entsteht, wenn zwei Transaktionen exklusive Rechte sperren, die die jeweils andere Transaktion bräuchte.
PostgreSQL erkennt diese automatisch und löst sie, in dem es eine Transaktion abbricht und dadurch die andere
erfolgreich abschließen lässt.
Die abgebrochene Transaktion gibt den Fehler \texttt{40P01 Deadlock detected} zurück~\cite{PostgreSQLDeadlocks}.

\texttt{55P03 lock not available} entsteht, wenn eine Transaktion versucht auf eine Zeile ein Lock zu erhalten,
welches nicht verfügbar ist, da dieses von einer anderen Transaktion schon benutzt wird~\cite{PostgreSQLRowLocks}.
Kann die Transaktion das Lock nicht innerhalb der konfigurierten Wartezeit erhalten, welche durch
\texttt{lock\_timeout} festgelegt wird, dann bricht PostgreSQL die Anfrage ab~\cite{PostgreSQLClientDefaults}.
Da die Transaktion explizit ein Lock mit \texttt{SELECT ... FOR UPDATE} anfordert und dieses Lock aufgrund des Timeouts
nicht erhält, wird die Transaktion mit dem Code \texttt{55P03 lock not available} abgebrochen.

Alle diese Fehler haben eine Sache gemeinsam, sie sind sogenannte \emph{transient Failures}.
Das bedeutet, dass sich die Probleme der Transaktionen von selbst wieder lösen.
Jeder Fehler der, oben erwähnt wurde, kann sich durch langes Warten lösen.
Wenn die abgebrochene Transaktion wartet, bis die andere erfolgreich abgeschlossen ist, dann kann bei nur zwei
Transaktionen kein Deadlock, Serialization Failure oder Lock not available mehr entstehen.
Deswegen sind diese Fehler transient Failures und sollten wiederholt werden~\cite{MicrosoftRetryPattern}.

Eine \texttt{Constraint-Violation} wäre kein transient Failure, da sich der Fehler nicht so schnell automatisch lösen
wird.